\begin{center}
\begin{tikzpicture}[>=stealth]
 \draw[->] (0,0) -- (8.5,0) node[below] {$x$};
 \draw[->] (0,0) -- (0,5) node[left] {$y$};
 \draw[->, thick, blue] (0,0) -- (6.71,0) node[midway, below] {$v_x$};
 \draw[->, thick, blue] (6.71,0) -- (6.71,4.36) node[midway, right] {$v_y$};
 \draw[->, very thick, red] (0,0) -- (6.71,4.36) node[midway, above left] {$\vec v_0$};
 \draw (1.6,0) arc (0:33:1.6);
 \node at (2.0,0.55) {$\alpha$};
\end{tikzpicture}
\end{center}

Der Geschwindigkeitsvektor $\vec v_0$ ist die Hypotenuse eines rechtwinkligen Dreiecks mit den Katheten $v_x$ (parallel zur Erdoberfläche) und $v_y$ (senkrecht dazu):
\begin{align}
 v_x &= \vxF \\
   &= \vab\cdot\cos\left(\alO\right) \\
   &= \vx \approx \result{\vxP{0}{2}} \\
 v_y &= \vyF \\
   &= \vab\cdot\sin\left(\alO\right) \\
   &= \vy \approx \result{\vyP{0}{2}}
\end{align}
